\chapter{哈尔测度　卷积}

希腊人的长度、面积和体积诸概念，本质上都建立在它们在位移下的不变性之上："彼此重合的事物（ἐφαρμόζοντα）相等"（欧几里得《几何原本》第一卷"共同观念"4）；正是通过巧妙地运用这条原理，人们得到了给出经典"图形"（多边形、圆锥曲线、多面体、球面等）面积或体积的全部公式，有时用有限分解的程序，有时用"穷竭法"。¹ 用现代语言来说，希腊几何学家所做的工作，可以视为证明了"集合函数"的存在性——这些函数可加且在位移下不变，但只对十分特殊类型的集合有定义。积分学可以看作是对扩大这类集合函数定义域这一需要的回应，而从卡瓦列里到 H. 勒贝格，正是这一关切处于分析学家研究的首位；至于位移下的不变性，它则沦为次要的考虑，因为这一性质已变成二重或三重积分中变量替换的一般公式以及正交变换行列式等于 ±1 这一事实的平凡推论。即便在非欧几何中（尽管那里的位移群有所不同），观点仍然相同：一般地，黎曼从 \(ds^2\) 出发，用经典的欧氏公式定义面积或体积的无穷小元素（或维数 ≥ 3 时的类似物），而它们在保持 \(ds^2\) 不变的变换下的不变性因此几乎是同义反复。

只是在 1890 年前后，群作用下不变测度概念的其他一些不那么直接的推广，才随着积分不变量理论——尤其是 H. 庞加莱和 E. 嘉当的工作——的发展而出现；

H. 庞加莱只预见到作用于空间一部分上的单参数群，而 E. 嘉当主要关注的是位移群，但它们作用的空间并非其定义所在的空间。例如，他由此（除其他结果外）确定了（{[}52 a{]}, v. II \(_{1}\), pp.~265-302）R \(^{2}\) 或 R \(^{3}\) 中直线空间上的（位移群下不变的）测度；\(^{2}\) 他还指出，李群的积分不变量一般说来无非是特殊的微分不变量，因而有可能用李的方法把它们全部确定出来。然而，在 1897 年 A. 胡尔维茨的那项基本工作{[}168{]}之前，似乎没有人想到在群自身上考虑或使用不变测度。胡尔维茨想要构造（R \(^{n}\) 上的）在正交群下不变的多项式，他从一个评注出发：对有限的线性变换群而言，只要对任意多项式 P 经群中所有元素 s 变换所得的象 s.P 取平均，问题便立即解决；这使他想到，对正交群，可以用关于一个不变测度的积分来代替取平均；他借助欧拉角的参数化表示明确给出了后者的表达式，但随即（独立于 E. 嘉当）注意到，李的方法可以给出所有李群上不变测度的存在性。也许正是由于二十世纪初不变量理论的衰落，胡尔维茨的思想几乎没有引起直接反响，直到 1924 年之后，随着 I. 舒尔和 H. 外尔把弗罗贝尼乌斯关于有限群线性表示的经典理论（第~121 页）推广到紧李群，这些思想才被认为有价值。前者仅限于正交群的情形，并说明胡尔维茨的方法如何能把特征标的经典正交关系推广开来；H. 外尔把这一思想与 E. 嘉当关于半单李代数的工作结合起来，得到了紧李群不可约表示特征标的显式表达式以及完全可约性定理{[}331 b{]}，随后又通过对"正则表示"概念的大胆推广，得到了著名的彼得–外尔定理，它与有限群理论中正则表示分解为其不可约分量的分解完美地类似{[}332{]}。

一年前，O. 施赖尔已经建立了拓扑群的一般理论{[}276 a{]}，从那时起人们已经清楚，彼得–外尔论文中的论证对于任何能在其上定义"不变测度"的拓扑群都将原样成立。说实话，当时拓扑学与测度论的一般概念尚处在全面形成之中，人们有望在其上定义不变测度的是哪一类拓扑群、这种"测度"又要对哪些集合来定义，看来都没有明确的界限。唯一明显的一点是：不能指望把证明李群上不变测度存在性的无穷小方法推广到一般情形。然而，另一股源于勒贝格测度方面工作的思想潮流，恰好引向更直接的求解方法。豪斯多夫在 1914 年证明，在 \(\mathbf{R}^3\) 的全部子集上有定义、在位移下不变且不恒为零的可加集合函数并不存在，于是自然要问这一结果对 \(\mathbf{R}\) 和 \(\mathbf{R}^2\) 是否仍然成立：S. 巴拿赫于 1923 年以一种出人意料的方式解决了这个问题，他证明恰恰相反，在这两种情形这样的"测度"确实存在{[}14 b{]}；他的做法十分巧妙，已经依赖于一个超限归纳构造，并用到函数被群中元素平移后的"平均"\(\frac{1}{n} \sum_{k=1}^{n} f(x + \alpha_k)\)。\(^3\) 正是这类思想使 A. 哈尔得以在 1933 年{[}139{]}迈出决定性的一步：证明具有可数开集基的局部紧群上不变测度的存在性。他受到经典积分学中用棱长任意小的全等立方体的拼接来逼近体积的方法的启发，借助对角线程序，把不变测度作为一列"逼近测度"的极限而得到，这一程序本质上至今仍在使用。这一发现影响巨大，尤其因为它使 J. 冯·诺伊曼能立即对紧群解决希尔伯特著名的"第五问题"，即只凭纯拓扑性质（排除一切预先给定的微分结构）刻画李群的问题。但人们随即注意到，要能有效地使用不变测度，仅知道它存在是不够的，还必须知道它在相差一个常数因子的意义下是唯一的；这一点首先由 J. 冯·诺伊曼对紧群证明，他用的是凭借连续函数的"平均"来定义哈尔测度的方法，与巴拿赫的方法类似{[}324 e{]}；随后 J. 冯·诺伊曼{[}324 f{]}与 A. 韦伊{[}330 c{]}用不同的方法同时得到了局部紧群情形的唯一性，A. 韦伊还同时指明了如何把哈尔的程序推广到一般的局部紧群。也是 A. 韦伊（出处同上）得到了齐性空间上相对不变测度的存在性条件，并最终证明，分离拓扑群上一个"测度"（具有合理的性质）的存在本身就意味着该群是局部预紧的。这些工作实质上完成了哈尔测度的一般理论；此后唯一值得引述的补充是拟不变测度的概念，它直到 1950 年前后才真正出现，与局部紧群在希尔伯特空间中的表示理论相关联。

卷积乘积的历史更为复杂。早在十九世纪初人们就注意到：例如，若 \(F(x,t)\) 是关于 x 与 t 的线性常系数偏微分方程的一个积分，则

\[
\int_ {- \infty} ^ {+ \infty} F (x - s, t) f (s) d s
\]

也是同一方程的积分；泊松等人在 1820 年之前就已经利用这一想法，把热方程的积分写成

\[
\int_ {- \infty} ^ {+ \infty} \exp \left(- \frac {(x - s) ^ {2}}{4 t}\right) f (s) d s.\tag{23.1}
\]

稍后，表达式

\[
\frac {1}{2 \pi} \int_ {- \pi} ^ {+ \pi} \frac {\sin \frac {2 n + 1}{2} (x - t)}{\sin \frac {x - t}{2}} f (t) d t\tag{23.2}
\]

给出傅里叶级数的部分和，而狄利克雷对该积分在 n 趋于 +∞ 时极限的研究，给出了环面 T 上"正则化" \(f \mapsto \rho_{n} * f\) 的第一个例子（实际上是由一列非正"核" \(\rho_{n}\) 作出的，这使它的研究大为复杂）；在"奇异积分"的名义下，类似的积分表达式将成为十九世纪末至二十世纪初的分析学家——从 P. 杜布瓦–雷蒙到 H. 勒贝格——所偏爱的课题。在 R 上，魏尔斯特拉斯利用积分 (23.1) 证明他那条关于多项式逼近的定理，并在此背景下给出了用正"核"序列 \(\rho_{n}\)（形如 \(x \mapsto c_{n} \rho(nx)\)）作正则化的一般原理。在 T 上，用正核作正则化的最著名例子稍后由费耶尔给出，从那时起，这一标准程序将成为绝大多数函数级数"求和法"的基础。

然而，这些工作由于"核"与被正则化的函数二者所扮演的角色缺少对称性，几乎没有显示出卷积乘积的代数性质。对这一点的强调尤其应归功于沃尔泰拉。他一般地研究了两个二元函数的"复合" \(F * G\)

\[
(F * G) (x, y) = \int_ {x} ^ {y} F (x, t) G (t, y) d t
\]

他将其视为两个矩阵的乘积经"从有限过渡到无穷"而得到的推广{[}322 a{]}。他很早便区分出（因其遗传理论中的解释而如此命名的）"闭循环"情形，其中 F 与 G 只依赖于 y - x；此时 \(H = F * G\) 亦然，而且若设 \(F(x, y) = f(y - x)\)，\(G(x, y) = g(y - x)\)，则有

\[
H (x, y) = h (y - x),
\]

其中

\[
h (t) = \int_ {0} ^ {t} f (t - s) g (s) d s
\]

于是，当 \(t \geq 0\) 时，\(h\) 与如下两个函数 \(f_1, g_1\) 的卷积相合：它们分别等于 \(f\) 与 \(g\)（对 \(t \geq 0\)），而对 \(t < 0\) 则等于 0。

与此同时，沃尔泰拉发展起来的代数形式体系并没有揭示出与 R 的群结构以及傅里叶变换的联系。我们在此无须撰写后者的历史；但应当指出，从柯西开始，研究傅里叶积分的分析学家主要关心的是为各种"反演"公式的有效性寻找越来越宽的条件，而多少忽略了它的代数性质。傅里叶本人的工作（以及拉普拉斯关于类似积分 \(\int_{0}^{+\infty} e^{-st} f(t) dt\) 的工作）当然不能这么说；但这些变换本质上是在与线性问题的关联中引入的，因此长久以来人们认为不值得考虑两个傅里叶变换的乘积也就不足为怪了（三角级数或整幂级数的乘积是例外，但它们与离散测度卷积的联系在十九世纪显然不可能被察觉）。这一乘积以及 R 上卷积的最早提及，大概见于切比雪夫的一篇论文{[}306{]}，与概率计算方面的问题有关。事实上，在概率论中，R 上两个"概率律"（总质量为 1 的正测度）的卷积 \(\mu * \nu\) 不是别的，正是 \(\mu\) 与 \(\nu\) 的"乘积"概率律（就相应"随机变量"之和而言）。当然，在切比雪夫那里，涉及的还只是具有（关于勒贝格测度的）密度的概率律的卷积，因而是函数的卷积；况且它在他那里只是偶一出现，在 1920–1930 年之前出现它的那少数几部著作中情形也是如此。1920 年，P.-J. 丹尼尔在一篇很少有人引用的短文{[}76 c{]}中定义了 R 上两个任意测度的卷积以及这种测度的傅里叶变换，并明确指出傅里叶变换把卷积变为通常的乘积；这一形式体系从 1925 年起将被概率学家们大量使用，尤其是追随 P. 莱维。但卷积在群论中的根本重要性直到 1927 年才由 H. 外尔充分认识到；他意识到，对紧群而言，函数的卷积扮演着有限群群代数中乘法的角色，从而使人能据此定义"正则表示"；同时，他还在正则化中找到了有限群代数中单位元的对应物。剩下要做的是把所有这些观点加以综合，这在 A. 韦伊的书中得以完成{[}330 c{]}，它是此后各种推广的前奏：一方面是 I. 盖尔范德的赋范代数，另一方面是广义函数的卷积。

哈尔测度与卷积很快成为不可或缺的工具，体现了如此鲜明地刻画着现代分析的代数化趋向；在后续的历史注记中，我们还须多次提到它们的应用。这里我们只限于指出其中涉及局部紧群闭子群（特别是离散子群）的"变化"的那一项应用。这一理论从数的几何中 K. 马勒的一个结果出发，由 C. 沙博蒂于 1950 年开创，最近则刚被 Macbeath 与 Swierczkowski 相当大地发展和深化了{[}212{]}。

